\model{Map and Set}
  \begin{center}
    \small
    \begin{tabular}{p{3.9in} p{1.9in}}
      \begin{minipage}{3.9in}
        \begin{cpplst}
#include <map>

map<string, int> population;
population["Seattle"] = 750000;
population["Bellingham"] = 90000;
population["Olympia"] = 52000;
population["Seattle"] += 1000;

for (pair<string,int> city : population) {
    cout << city.first << ": ";
    cout << city.second << endl;
}
cout << population.count("Seattle");
cout << population.count("Portland");
        \end{cpplst}
      \end{minipage}
      &
      \begin{minipage}{1.9in}
        {\bf population map:}\\[4pt]
        \renewcommand{\arraystretch}{1.2}
        \begin{tabular}{|c|c|}
          \hline
          \rowcolor{orange!20} \textbf{key} & \textbf{value} \\
          \hline
          \tt "Bellingham" & 90000 \\
          \hline
          \tt "Olympia"    & 52000 \\
          \hline
          \tt "Seattle"    & 751000 \\
          \hline
        \end{tabular}\\[10pt]
        {\bf Output:}
        \hrule\vskip 5pt
        {\tt Bellingham: 90000}\\
        {\tt Olympia: 52000}\\
        {\tt Seattle: 751000}\\
        {\tt 1}\\
        {\tt 0}
      \end{minipage}
    \end{tabular}
    \vskip 10pt
    \begin{tabular}{p{3.9in} p{1.9in}}
      \begin{minipage}{3.9in}
        \begin{cpplst}
#include <set>

set<int> scores = {85, 92, 78, 92, 85, 100};
scores.insert(73);
scores.insert(92);

for (int score : scores) {
    cout << score << endl;
}
cout << "size: " << scores.size();
        \end{cpplst}
      \end{minipage}
      &
      \begin{minipage}{1.9in}
        {\bf scores set (sorted, unique):}\\[4pt]
        \begin{tabular}{|c|c|c|c|c|}
          \hline
          73 & 78 & 85 & 92 & 100 \\
          \hline
        \end{tabular}\\[10pt]
        {\bf Output:}
        \hrule\vskip 5pt
        {\tt 73}\\
        {\tt 78}\\
        {\tt 85}\\
        {\tt 92}\\
        {\tt 100}\\
        {\tt size: 5}
      \end{minipage}
    \end{tabular}
  \end{center}

  {\it\large Refer to Model 3 above as your team develops consensus answers
    to the questions below.}

  \quest{20 min}

  \Q A {\tt map} stores key-value pairs, where each key maps to exactly
    one value.  Answer the following about the {\tt map} example above.
    \begin{enumerate}
      \itemsep 10pt
      \item What is the type of the {\it key} in {\tt population}?
        \hfill\ans[1.5in]{\cpp{string}}

      \item What is the type of the {\it value} in {\tt population}?
        \hfill\ans[1.5in]{\cpp{int}}

      \item What syntax is used to insert or update a value in the map?
        \hfill\ans[2in]{\cpp{map[key] = value}}

      \newpage

      \item What happens when {\tt population["Seattle"]} is assigned
        a second time?  Does it create a new entry or modify the existing one?

        \vskip -5pt

        \begin{answer}[0.5in]
          It modifies the existing entry.  A map key is unique; assigning
          to an existing key overwrites its associated value.
        \end{answer}
    \end{enumerate}

  \vskip -20pt

  \Q What does {\tt map.count(key)} return?  What are the two possible
    return values and what does each mean?
    \begin{answer}[0.5in]
      It returns {\tt 1} if the key exists in the map, and {\tt 0} if
      it does not.  (A map can hold at most one entry per key.)
    \end{answer}


  \Q When iterating over a {\tt map} or {\tt set} with a range-based
    for loop, the elements are visited in a particular order.\key\\[-2mm]
    \begin{enumerate}
      \itemsep 10pt
      \item In what order does the {\tt for} loop print the cities from
        the {\tt map}?
        \hfill\ans[1.5in]{Alphabetical (sorted by key)}

      \item In what order does the {\tt for} loop print the values from
        the {\tt set}?
        \hfill\ans[1.5in]{Ascending numeric order}

      \item Both {\tt map} and {\tt set} keep their elements
        {\it sorted}.  How does this differ from {\tt vector} and {\tt list}?

        \vskip -5pt

        \begin{answer}[0.5in]
          A \cpp{vector} and \cpp{list} maintain {\it insertion order};
          elements appear in the order you added them.  A \cpp{map}
          and \cpp{set} always store elements in sorted order
          regardless of insertion order.
        \end{answer}
    \end{enumerate}

  \vskip -20pt

  \Q The {\tt set} in the model is initialized with six values, two
    of which are duplicates (92 appears twice, 85 appears twice).
    \begin{enumerate}
      \itemsep 10pt
      \item How many elements does the set actually contain?
        \hfill\ans[0.75in]{5}

      \item After {\tt scores.insert(92)}, does the set size increase?
        \hfill\ans[0.75in]{No}

      \item When would a {\tt set} be more useful than a {\tt vector}
        for storing a collection of values?

        \vskip -5pt

        \begin{answer}[0.5in]
          A \cpp{set} is useful when you need to store only unique
          values, automatically eliminate duplicates, or quickly test
          whether a value is present (using \cpp{count}).
        \end{answer}
    \end{enumerate}

  \vskip -20pt

  \Q The file {\tt activity28c.cpp} contains both examples.  Run it
    and verify the output.  Then write C++ code to count the number
    of times each word appears in the sentence {\tt "the cat sat on
    the mat"}.  Use a {\tt map<string,int>}.
    \begin{answer}[1.05in]
      \begin{cpplst}
map<string, int> wordCount;
vector<string> words =
    {"the","cat","sat","on","the","mat"};
for (string word : words) {
    wordCount[word]++;
}
for (pair<string,int> entry : wordCount) {
    cout << entry.first << ": ";
    cout << entry.second << endl;
}
      \end{cpplst}
    \end{answer}
